\documentclass[12pt]{article}
\usepackage[utf8]{inputenc}
\usepackage[T1]{fontenc}
\usepackage{hyperref}
\title{Sumerian Language Resources for K-12 Curriculum Development}
\author{Compiled by Claude Sonnet 5.5 (claude-sonnet-5-5) for Andr\'as Kornai}
\date{}
\begin{document}
\maketitle

\section{Introduction}
Sumerian (ISO 639-3 sux; Glottocode sume1241) is an ancient language isolate of Mesopotamia with no native speakers today. Any K-12 programme would be a heritage and classical-language programme. It would use Sumerian for as much as is feasible, for example sign reading, transliteration and short texts. The resources below are the ones found in a web search. They are mostly scholarly and in English.

\section{Primary Educational Resources}
No school curriculum for Sumerian was found. The following are university-level teaching grammars and readers that could be adapted.
\begin{itemize}
\item Foxvog, \emph{Introduction to Sumerian Grammar}: free online via CDLI; includes a grammar, a writing-system overview and exercises. Reported in the search results; no direct URL was returned, so none is given here.
\item Getting Started with Sumerian (ISAW Library, NYU), a guide to beginning study: \url{https://isaw.nyu.edu/library/blog/getting-started-with-sumerian}
\end{itemize}

\section{Linguistic and Reference Materials}
\begin{itemize}
\item Hayes, \emph{Manual of Sumerian Grammar and Texts}, Zólyomi, \emph{An Introduction to the Grammar of Sumerian}, Jagersma, \emph{A Descriptive Grammar of Sumerian}, and Volk, \emph{A Sumerian Chrestomathy} are named in the search results as standard works. No free online URLs were found for them.
\item Flashcards for Hayes's manual: \url{https://mnemosyne-proj.org/cards/hayess-manual-sumerian-grammar-and-texts}
\item ePSD2, the electronic Pennsylvania Sumerian Dictionary: \url{https://stars.library.ucf.edu/ancientneareast/26}
\end{itemize}

\section{Religious and Community Resources}
No community organisations or living religious communities were found. Sumerian literature (hymns, myths, proverbs) survives in cuneiform texts. A survey of online resources for ancient Mesopotamia: \url{https://digitalorientalist.com/2019/03/06/online-resources-for-ancient-mesopotamia/}

\section{Digital and Audio Resources}
ePSD allows searching by individual cuneiform sign, which helps beginners. No audio or video resources were found in the search results. Pronunciation is reconstructed.

\section{Curriculum Development Context}
There are no native speakers. Little material exists at school level, and what exists is university-oriented. A K-12 course would need adapted materials built on the grammars and the ePSD lexicon, with expert review by Assyriologists. Verify every resource before use.
\end{document}